% Part 2: System Model and Problem Formulation

\section{System Model and Problem Formulation}

\subsection{LEO Satellite Network Model}

We consider a LEO satellite constellation $\mathcal{S} = \{s_1, s_2, ..., s_N\}$ consisting of $N$ satellites distributed across multiple orbital planes. The constellation topology at time $t$ is represented as a directed graph $\mathcal{G}(t) = (\mathcal{V}, \mathcal{E}(t))$, where:

\begin{itemize}
\item $\mathcal{V} = \mathcal{S}$ is the set of vertices (satellites)
\item $\mathcal{E}(t) \subseteq \mathcal{V} \times \mathcal{V}$ is the set of inter-satellite links (ISLs) at time $t$
\end{itemize}

Each satellite $s_i$ is characterized by its position vector $\mathbf{p}_i(t) = [x_i(t), y_i(t), z_i(t)]$ in Earth-Centered Inertial (ECI) coordinates, which evolves according to:

\begin{equation}
\mathbf{p}_i(t + \Delta t) = \mathbf{p}_i(t) + \mathbf{v}_i(t) \cdot \Delta t
\end{equation}

where $\mathbf{v}_i(t)$ is the velocity vector and $\Delta t$ is the simulation timestep.

An ISL exists between satellites $s_i$ and $s_j$ if:

\begin{equation}
d_{ij}(t) = \|\mathbf{p}_i(t) - \mathbf{p}_j(t)\| \leq d_{\max}
\end{equation}

where $d_{ij}(t)$ is the Euclidean distance and $d_{\max}$ is the maximum ISL range. We also enforce line-of-sight constraints to prevent links that pass through Earth.

\subsection{Link Characterization}

Each ISL $e_{ij} \in \mathcal{E}(t)$ is characterized by:

\begin{itemize}
\item \textbf{Propagation Delay}: $\tau_{ij}(t) = \frac{d_{ij}(t)}{c}$, where $c$ is the speed of light
\item \textbf{Capacity}: $C_{ij}$ (Mbps) - the maximum data rate
\item \textbf{Queue Length}: $Q_{ij}(t)$ - number of packets buffered
\item \textbf{Utilization}: $U_{ij}(t) = \frac{B_{ij}(t)}{C_{ij}} \times 100\%$, where $B_{ij}(t)$ is the instantaneous bandwidth usage
\end{itemize}

The total end-to-end delay for a path $\mathcal{P} = (s_1, s_2, ..., s_k)$ is:

\begin{equation}
D_{\mathcal{P}}(t) = \sum_{(s_i,s_{i+1}) \in \mathcal{P}} \left[\tau_{i,i+1}(t) + d_{\text{proc}} + Q_{i,i+1}(t) \cdot d_{\text{queue}}\right]
\end{equation}

where $d_{\text{proc}}$ is the processing delay per hop and $d_{\text{queue}}$ is the per-packet queuing delay.

\subsection{Traffic Model}

We model three traffic slices corresponding to 6G service classes:

\begin{enumerate}
\item \textbf{URLLC (Ultra-Reliable Low-Latency Communications)}: Critical applications requiring latency $< 10$ ms and reliability $> 99.999\%$. Examples include autonomous vehicles and industrial automation.

\item \textbf{eMBB (Enhanced Mobile Broadband)}: High-throughput applications with moderate latency tolerance ($< 50$ ms). Examples include video streaming and cloud gaming.

\item \textbf{mMTC (Massive Machine-Type Communications)}: IoT applications with relaxed latency requirements ($< 200$ ms) but massive connectivity needs.
\end{enumerate}

Each traffic slice $k \in \{\text{URLLC}, \text{eMBB}, \text{mMTC}\}$ is characterized by:

\begin{equation}
\text{Slice}_k = \{\lambda_k, L_k^{\max}, \pi_k, w_k\}
\end{equation}

where:
\begin{itemize}
\item $\lambda_k$ is the packet arrival rate
\item $L_k^{\max}$ is the maximum tolerable latency
\item $\pi_k$ is the priority level (1 = highest)
\item $w_k$ is the traffic weight/proportion
\end{itemize}

\subsection{Problem Formulation}

Given the dynamic network graph $\mathcal{G}(t)$ and traffic demands, our objective is to learn a routing policy $\pi_\theta$ parameterized by $\theta$ that maximizes the expected cumulative reward:

\begin{equation}
\max_{\theta} \mathbb{E}_{\pi_\theta}\left[\sum_{t=0}^{T} \gamma^t R_t\right]
\end{equation}

where $\gamma \in [0,1]$ is the discount factor and $R_t$ is the reward at timestep $t$.

The reward function is designed to balance multiple objectives:

\begin{equation}
\begin{split}
R_t = & \alpha_1 \cdot R_{\text{delivery}}(t) + \alpha_2 \cdot R_{\text{QoS}}(t) \\
      & - \beta_1 \cdot P_{\text{latency}}(t) - \beta_2 \cdot P_{\text{drops}}(t) \\
      & - \beta_3 \cdot P_{\text{congestion}}(t)
\end{split}
\end{equation}

where:

\begin{itemize}
\item $R_{\text{delivery}}(t) = \mathbb{1}_{\text{delivered}} \times 100$ rewards successful packet delivery
\item $R_{\text{QoS}}(t) = \mathbb{1}_{\text{QoS met}} \times 50$ rewards meeting QoS targets
\item $P_{\text{latency}}(t) = \frac{D_t}{L_k^{\max}}$ penalizes high latency relative to target
\item $P_{\text{drops}}(t) = 500$ heavily penalizes packet drops
\item $P_{\text{congestion}}(t) = \sum_{e_{ij}} \frac{Q_{ij}(t)}{Q_{\max}}$ penalizes network congestion
\end{itemize}

The hyperparameters $\alpha_1, \alpha_2, \beta_1, \beta_2, \beta_3$ control the trade-offs between objectives.

\subsection{State Space}

At each routing decision, the agent observes state $\mathbf{s}_t$ consisting of:

\begin{equation}
\mathbf{s}_t = [\mathbf{s}_t^{\text{global}}, \mathbf{s}_t^{\text{neighbor}}]
\end{equation}

\textbf{Global Features} (16 dimensions):
\begin{itemize}
\item Current satellite position (3D coordinates)
\item Destination satellite position
\item Packet priority and slice type
\item Current queue depth
\item Hop count and time-to-live
\item Geographic progress toward destination
\end{itemize}

\textbf{Per-Neighbor Features} (13 dimensions $\times$ $N_{\max}$ neighbors):

For each neighboring satellite $j$:
\begin{itemize}
\item Propagation delay $\tau_{ij}$
\item Queue length $Q_{ij}$
\item Link utilization $U_{ij}$
\item Bandwidth availability
\item Success rate (historical)
\item Packet drop rate
\item Distance to destination
\item Angular velocity
\item Link stability indicator
\item Congestion level
\item Average delay
\item \textbf{Predicted queue length} $\hat{Q}_{ij}(t+h)$ (NEW)
\item \textbf{Predicted utilization} $\hat{U}_{ij}(t+h)$ (NEW)
\end{itemize}

The predicted features enable proactive routing by anticipating future link conditions.

\subsection{Action Space}

The action space consists of selecting the next-hop satellite from the set of available neighbors:

\begin{equation}
\mathcal{A}_t = \{s_j : e_{ij} \in \mathcal{E}(t), j \neq i\}
\end{equation}

The action is represented as a discrete choice $a_t \in \{0, 1, ..., N_{\max}-1\}$, where $N_{\max}$ is the maximum number of neighbors (padded if necessary).

\subsection{Challenges and Design Considerations}

The problem presents several challenges:

\begin{enumerate}
\item \textbf{Partial Observability}: The agent only observes local neighbor information, not global network state.

\item \textbf{Non-Stationary Environment}: Satellite motion causes continuous topology evolution.

\item \textbf{Multi-Objective Optimization}: Balancing latency, throughput, reliability, and fairness requires careful reward shaping.

\item \textbf{Scalability}: The action space grows with the number of neighbors, and the network contains hundreds of satellites.

\item \textbf{Cold Start}: Prediction models require historical data, creating a warmup period.
\end{enumerate}

Our TA-PPO framework addresses these challenges through:
\begin{itemize}
\item Decentralized decision-making at each satellite
\item Continuous learning and adaptation
\item Hierarchical reward structure
\item Efficient neural network architectures
\item Adaptive prediction warmup strategy
\end{itemize}

